\section{Evaluation and Reproducibility}
\label{sec:evaluation}

\subsection{Recorded execution costs}

The fixed-host certificate-path profile uses one warmup and seven measured
runs for each stage on a \(2\times3\) integer SNF input. Table
\ref{tab:certificate-costs} reports medians and interquartile ranges from the
committed JSON and CSV observations. The strict kernel stages include
elaboration and kernel checking of a generated certificate module; they are
not comparable to the native Boolean checker as interchangeable trust modes.

\begin{table}[ht]
  \centering
  \small
  \begin{tabular}{lrrr}
    \toprule
    Stage & Median wall (s) & IQR (s) & Median RSS (KiB) \\
    \midrule
    Internal Lean generation & 0.10 & 0.00 & 64,304 \\
    FLINT process generation & 0.03 & 0.00 & 20,836 \\
    FLINT FFI generation & 0.06 & 0.00 & 66,744 \\
    Strict kernel checking & 7.04 & 0.07 & 4,178,000 \\
    Native Boolean checking & 0.02 & 0.00 & 64,464 \\
    Process generation and import & 7.33 & 0.27 & 4,144,616 \\
    FFI generation and import & 7.30 & 0.16 & 4,171,756 \\
    \bottomrule
  \end{tabular}
  \caption{Certificate-path observations on the recorded Ryzen 9 7940H host.}
  \label{tab:certificate-costs}
\end{table}

Both generated JSON certificates occupy 1,149 bytes and the emitted Lean
module occupies 3,888 bytes. These figures describe one small input and one
hardware fingerprint. Hosted CI does not use them as absolute time limits.
Automatic comparisons require an identical fingerprint; otherwise a run
establishes a new observational baseline. Ring-operation, pivot, and gcd
counts remain explicit null fields because neither current execution path
exposes truthful counters.

The separate rational-polynomial native corpus covers pivot improvement, zero
and rectangular matrices, rank deficiency, transformation equations, and all
stored inverse identities. Its seven-run median is \(4.66\) seconds with a
\(1.29\)-second IQR and 72,056 KiB median maximum RSS on the same host.

\subsection{Trust boundary}

The theorem-producing path is intentionally narrower than the execution
pipeline. The trusted conclusion is a Lean term checked by the kernel against
the declared axioms. The executable \code{Fin} kernels and strict imported
certificate declarations are audited declaration by declaration; both tiers
contain only \code{propext} and \code{Quot.sound}. The arbitrary finite-index
convenience layer may additionally use \code{Classical.choice}, and its
separate facade-wide report records that fact.

The \code{decide_cbv} tactic, parser, source emitter, C worker, Python wrapper,
FLINT, FFI bridge, and native Boolean checker may all fail or contain bugs
without directly extending the kernel's axiom set. They are responsible for
producing source or proof terms that the kernel checks. In particular,
successful native or differential execution is never reported as a theorem.
The strict source modules reject \code{native_decide},
\code{Lean.ofReduceBool}, and \code{Lean.ofReduceNat}.

\subsection{Artifact profile}

The core artifact pins Lean 4.32.0 by archive hash, mathlib and every Lake
dependency by \code{lake-manifest.json}, and the Debian base and Dockerfile
frontend by OCI digest. Its release gate rebuilds the library, exact API
freeze, tests, deterministic execution corpus, three axiom reports, generated
certificate modules, standalone mathlib patch, metadata, and manuscript. The
optional FLINT profile additionally pins FLINT 3.6.0, GMP 6.3.0, and MPFR
4.2.2 by source hash and ships the corresponding sources, licenses, build
configuration, and relink instructions.

The archive builder operates only on a clean committed tree. It emits the
source archive, manuscript, citation/deposit metadata, build information, an
optional FLINT source bundle, and a checked SHA256 manifest. It performs no
publication, tagging, deposition, or upstream submission.
